\documentclass[letterpaper]{article} % DO NOT CHANGE THIS
\usepackage[submission]{aaai2027}  % DO NOT CHANGE THIS
\usepackage[hyphens]{url}  % DO NOT CHANGE THIS
\usepackage{graphicx} % DO NOT CHANGE THIS
\usepackage{natbib}  % DO NOT CHANGE THIS
\usepackage{caption} % DO NOT CHANGE THIS
\usepackage{amsmath}
\usepackage{amssymb}
\usepackage{booktabs}
\frenchspacing  % DO NOT CHANGE THIS
\setlength{\pdfpagewidth}{8.5in} % DO NOT CHANGE THIS
\setlength{\pdfpageheight}{11in} % DO NOT CHANGE THIS

\newcommand{\method}[1]{\texttt{#1}}

% AAAI-27 compressed version of the Q/R/M/C measurement-framework
% manuscript.  The full single-column version (42 pp.) is the
% Supplementary Document; this file keeps 7 pages of technical content
% per the compression plan (docs/AAAI_COMPRESSION_PLAN.md, revised
% 2026-07-04: the external-framework validation moved INTO the body).

\pdfinfo{
/TemplateVersion (2027.1)
}

\title{Q/R/M/C: A Stage-Decomposition Evaluation Protocol\\for Memory-Based Agents}
\author{Anonymous submission}
\affiliations{Anonymous submission}

\begin{document}

\maketitle

\begin{abstract}
Success-only metrics collapse qualitatively different memory failures
into a single scalar and hide which stage of the retrieval pipeline is
responsible. We describe an evaluation protocol, the Q/R/M/C
decomposition (Query formation, Retrieval satisfaction, target
Materialization, Completion), operationalised so that four conditional
rates can be computed from standard trajectory logs. The arithmetic is
an elementary chain rule over logged events; the substance is in the
event definitions, which make the rates estimable in practice, much as
precision and recall are useful because of their definitions and not
their arithmetic. In single-agent experiments (MiniGrid, BabyAI),
oracle-stage interventions separate a retrieval-limited task
($0.39\!\to\!0.60$ under oracle retrieval) from downstream-limited
ones, and log inspection refines the diagnosis: ranking among
retrieved candidates is near-optimal ($96.9\%$), but the candidate set
is empty at $91\%$ of decision points, so the limiting factor is
memory accumulation. In multi-agent experiments (35{,}640 runs, four
sharing architectures, cadence $K$), the predicted bottleneck shift
between Materialization and Completion is confirmed as cluster-robust
slopes (Spearman $-0.53$/$+0.43$), time-to-success among successes
has an interior optimum at $K{=}8$, and a provenance annotation
identifies the mechanism: under fast sharing the completion collapse
is concentrated in locks backed by peers' evidence. The same tool-trace contract also
instruments three third-party agent stacks (OpenAI SDK, LangGraph,
AutoGen). Engineered structural failures localise identically on every
stack ($R^\star{=}0$, $C^\star{=}0$ in $20/20$), while on clean tasks
the frameworks spread $0.50$--$0.75$ end-to-end, and the stages
attribute the spread to weak first-commitment discipline combined with
different recovery loops. The protocol makes memory failures
comparable across symbolic and LLM systems, single- and multi-agent,
wherever a query/lock interface is exposed; on opaque learned
policies it measures whether the stages are observable at all.
\end{abstract}

\section{Introduction}

Embodied, goal-conditioned, and LLM agents increasingly retrieve
places, objects, and facts \emph{by meaning} from an explicit memory,
yet they are evaluated almost exclusively end-to-end: if success
improves, memory is judged useful. Two systems with the same success
rate can fail for different reasons. One never accumulates queryable
evidence; another retrieves perfectly and then loses the target to a
peer. This paper describes a measurement instrument that separates
such cases.

The instrument is a four-stage decomposition of memory-based task
completion: \textbf{Q}uery formation (the planner emits a semantic
intent), \textbf{R}etrieval satisfaction (memory returns a candidate
matching the task), target \textbf{M}aterialization (the planner locks
one candidate as an actionable target), and \textbf{C}ompletion (the
controller reaches it). Logged as per-episode events
$Q^\star,R^\star,M^\star,C^\star$, they factorise semantic-path
completion into conditional rates estimable from ordinary trajectory
logs (Definition~1; a multi-agent extension with memory and occupancy
coupling channels is Definition~2). The consistency argument is the
chain rule plus frequency estimation, and we claim no mathematical
novelty for it. The precedent is the decomposition of accuracy into
precision and recall, which was also arithmetically trivial and still
changed evaluation practice. Our contribution is a logging discipline
and the diagnoses it supports.

\paragraph{Contributions.}
\begin{enumerate}
\item \textbf{The protocol.} Operational, log-level definitions of the
  four stage events for single agents (Definition~1) and communicating
  collectives (Definition~2), with identifiability under a stated
  conditional stage-Markov assumption, plus a stress test of that
  assumption.
\item \textbf{Stage-level diagnoses.} Single-agent:
  oracle-stage interventions and log forensics attribute failures to
  memory accumulation rather than to ranking. Multi-agent: a
  registered slope-level claim (Empirical Claim~1: fast sharing
  raises $P(M^\star|R^\star)$ and lowers $P(C^\star|M^\star)$) is
  confirmed on a 35{,}640-run sweep, with an interior time-to-success
  optimum at $K{=}8$ and a provenance-level mechanism.
\item \textbf{External validity.} A two-layer contract (invariant
  event roles plus a system-specific adapter) instruments OpenAI-SDK,
  LangGraph, and AutoGen agents with one shared tool-trace adapter and
  no per-framework changes; engineered structural failures localise
  identically on all stacks, and behavioural differences between
  stacks are attributed stage-by-stage.
\end{enumerate}

\begin{table}[t]
\caption{Portability of the contract across system classes: the same
four events, defined once, instrument each class without modification
and localize a different failure in each. Sections
\ref{sec:single}--\ref{sec:external} instantiate the rows in order.}
\label{tab:portability}
\centering
\small
\setlength{\tabcolsep}{3pt}
\begin{tabular}{p{2.5cm} p{5.4cm}}
\toprule
System class & What the same four events revealed \\
\midrule
Symbolic single agent & ranking $96.9\%$ vs.\ empty candidates $91\%$: an \emph{accumulation} limit \\
Symbolic collective & M$\leftrightarrow$C bottleneck shift; $t_{\mathrm{succ}}$ optimum $K{=}8$; provenance mechanism \\
Learned policies & competitive success, collapsed stage observability \\
LLM agents & $1.00\!\to\!0.00$ collapse $=$ pure C-budget; sharing lowers success $2/3\!\to\!1/3$ \\
Third-party stacks $\times$ 2 models & structural failures localize identically; $M_1$ separates models; frameworks reshape it \\
\bottomrule
\end{tabular}
\end{table}

\noindent\fbox{\parbox{0.97\columnwidth}{\small
\textbf{Scope and non-claims.} Each limitation is stated together with
the reason the results survive it.
(i)~The contribution is the operationalisation: log-level event
definitions under which the four conditional rates are identifiable on
any system exposing a query/lock interface. The factorisation
arithmetic itself (the chain rule) carries no mathematical novelty,
and none is claimed.
(ii)~Empirical Claim~1 is a slope-level empirical observation rather
than a theorem, stated deliberately so that each slope is separately
falsifiable on logs; both were confirmed with cluster-robust
confidence intervals.
(iii)~No claim is made about the conditional product $\Phi$, which is
monotone in our data; the interior optimum lives on conditional
time-to-success, and we report the censoring-aware contrast under
which it does not persist.
The operationalisation is testable on specific summaries and can be
wrong on others; a protocol without such exposure would be
unfalsifiable.
(iv)~Primary substrates are grid worlds because they admit oracle
interventions and exhaustive sweeps (35{,}640 runs) that richer
substrates do not; the directional VMAS check and the LLM and
third-party studies indicate that the event definitions do not depend
on grid structure.
(v)~Multi-agent place identity assumes a shared coordinate frame in
the main sweep; a companion manuscript removes the assumption
constructively, leaving the event definitions unchanged.
(vi)~$R^\star$ and $M^\star$ presuppose an entity correspondence
between memory content and world referents, supplied here by
instrumented environments. The requirement is mild in the common
evaluation setting of this paper --- interchangeable agents (a shared
perception interface, possibly different weights) situated in a single
environment inherit the correspondence from the shared substrate ---
and becomes a genuinely open problem only for heterogeneous agents in
open worlds, where the protocol degrades gracefully to $Q^\star$ and
$C^\star$.
}}

\section{The Q/R/M/C Protocol}
\label{sec:protocol}

An episode issues a semantic query (required/preferred/penalty
tags); memory returns candidates; the planner locks one; the
controller executes. The four events anchor at logged triggers on the
episode's query--lock chain and are \emph{nested by construction}:
$Q^\star$ fires on intent emission; $R^\star := Q^\star \wedge {}$(the
returned candidate set contains a task-satisfying candidate);
$M^\star := R^\star \wedge {}$(the planner locks a true target drawn
from that candidate set); $C^\star := M^\star \wedge {}$(the
controller reaches the locked target). An episode may run several
query--retrieve--lock attempts; a link fires if some attempt realizes
it together with its predecessors, so $C^\star \subseteq M^\star
\subseteq R^\star \subseteq Q^\star$ holds at the episode level by
definition. $C^\star$ is \emph{semantic-path completion} ---
completion through the retrieval path --- and is distinct from raw
task success, which can also occur off-path.

We keep two kinds of quantities separate. \emph{Link rates} are the
conditionals $q = P(Q^\star)$, $r = P(R^\star \mid Q^\star)$,
$m = P(M^\star \mid R^\star)$, $c = P(C^\star \mid M^\star)$,
estimated by conditional frequency over episodes; under the nesting
above, $P(C^\star) = q\,r\,m\,c$ is an exact identity --- a chain
rule that needs no further assumption. \emph{Marginal event
frequencies} are used as four separate diagnostics on systems whose
logs expose no lock chain (\S\ref{sec:external}); there an event can
occur without its predecessor, the events are not nested, and no
product is claimed.

\textbf{Definition 1} (single-agent estimability, informal). With the
nested events above, $P(C^\star) = q\,r\,m\,c$ holds identically and
each link rate is consistently estimable from logs. The conditional
stage-Markov property (each stage outcome depends on the trace only
through its predecessor stage) is not needed for the identity; it is
what licenses reading each link as a stable property of the owning
subsystem, and hence predicting the effect of a stage intervention. A
simulation stress test quantifies what breaks when it fails: with a
hidden confounder, $m$ stays unbiased while $c$ inflates, so
diagnoses that hinge on the M--C \emph{contrast} degrade gracefully
(supplement).

\textbf{Definition 2} (multi-agent extension, informal). With $N$
agents sharing memory by peer broadcast every $K$ ticks, the
conditioning sets enlarge with two coupling channels: the shared
memory $\mathcal{M}^{(K)}_{-i}$ and occupancy $\mathcal{O}_{-i}$
(scarce targets are consumed). The same four rates are logged per
agent.

\textbf{Empirical Claim 1} (bottleneck shift, slope-level). As $K$
decreases (faster sharing), $P(M^\star|R^\star)$ rises, because
fresher peer evidence enables reliable commitment, and
$P(C^\star|M^\star)$ falls, because peers race for the same targets.
Both are claims about conditional rates, not about their product.

Interpretive readings of the stages (an order-theoretic account of
empty retrieval via the query--extent Galois connection; an
options/semi-MDP account of why optima appear on time-bearing metrics)
are confined to the supplement and are not counted among the
contributions.

\section{Single Agent: Localizing the Failure}
\label{sec:single}

On a 10-seed MiniGrid benchmark (water/rest semantic search,
goal-region rejoin in FourRooms, hazard recovery in LavaGapS7; BabyAI
portability in the supplement), end-to-end success hides
heterogeneity: water and rest are strong ($0.95$, $0.93$), while
rejoin and hazard recovery are weak. The reasons differ, and
oracle-stage interventions separate them. Replacing retrieval with an
oracle lifts hazard-recovery success $0.39\!\to\!0.60$ (bootstrap
95\% CIs $[0.35,0.43]$ vs.\ $[0.51,0.69]$; per-seed paired deltas
positive in $9/10$ seeds, exact sign test $p{=}0.004$;
retrieval-limited); the same oracle rescues semantic-path completion
on goal-region rejoin without changing end-to-end success
(downstream-limited). The framework's own logs refine the attribution
further: when candidates exist, ranking precision is $96.9\%$, but the
candidate set is empty at $91\%$ of decision points --- a figure
robust to the weighting ($0.908$ pooled over calls, $0.909$
seed-weighted, $0.891$ episode-weighted), so it is not an artefact of
long failing episodes contributing more decision points. The binding
constraint is accumulation, since the agent rarely has anything to
rank. In our own development this number redirected effort from
rankers to consolidation; success-only evaluation gives no indication
of it. Accumulation itself is not one of the four stages: the
protocol localises the failure to the R link, and within-link log
forensics names the cause. We therefore read Q/R/M/C as coarse
localisation with a designated drill-down, not as a complete causal
ontology of memory failure.

\begin{figure}[t]
  \centering
  \includegraphics[width=0.98\columnwidth]{figures/fig_oracle_waterfall.pdf}
  \caption{The diagnose$\to$intervene$\to$verify loop on the
  single-agent benchmark: oracle-stage interventions lift exactly the
  stage that the decomposition identifies. Hazard recovery is
  retrieval-limited (oracle-R: $0.39\!\to\!0.60$ end-to-end);
  goal-region rejoin is downstream-limited (oracle-R rescues
  semantic-path completion, end-to-end unchanged).}
  \label{fig:oracle}
\end{figure}

A learned behaviour-cloning baseline illustrates the observability
caveat: it reaches competitive success with collapsed stage structure
(there is no explicit query/lock interface), so end-to-end
optimisation removes exactly the events the protocol needs.

\section{Multi-Agent: Memory as a Shared Resource}
\label{sec:multi}

\begin{figure*}[t]
  \centering
  \includegraphics[width=0.92\textwidth]{figures/fig_conditional_rates.pdf}
  \caption{Per-architecture conditional Q/R/M/C rates (top) and mean
  time-to-success (bottom) on the 35{,}640-run sweep. Shared bus,
  central aggregator, and fast peer broadcast ($K{\le}4$) sit in the
  M-saturation/C-collapse regime; independent and slow peer
  ($K{\ge}16$) in the reverse. Error bars and band: 95\%
  cluster-bootstrap CIs over the 81 design cells. Time-to-success
  among successful runs has an interior minimum at $K{=}8$ (7.63 vs.\
  8.04 independent); the sharp statement is the cell-paired contrasts
  in the text, and the censoring-aware counterpart is also given
  there.}
  \label{fig:conditional-rates}
\end{figure*}

Four sharing architectures run a scarce water-search task over one
symbolic substrate: independent private memories, a shared event bus,
a central trust-weighted aggregator, and peer-to-peer broadcast at
cadence $K{\in}\{1,\dots,64\}$ ($N{\in}\{3,5,8\}$ agents, $T{\le}N$
targets, three layouts, three hazard densities, 40 seeds; 35{,}640
runs). Figure~\ref{fig:conditional-rates} shows the stage profile of
every architecture; retrieval saturates, and the differences between
architectures are confined to the M and C links, as Empirical Claim~1
predicts.

\textbf{Bottleneck shift.} Across peer runs, Spearman of
$P(M^\star|R^\star)$ vs.\ $K$ is $-0.53$ (cluster-robust 95\% CI
$[-0.59,-0.47]$ over 81 configuration cells) and of
$P(C^\star|M^\star)$ vs.\ $K$ is $+0.43$ ($[+0.37,+0.50]$). The
sharper signature is within-cadence: Pearson between the two links
inside a fixed $K$ peaks at $-0.61$ at $K{=}4$ and decays to noise by
$K{=}64$, so the bottleneck migrates between stages within the same
configuration as $K$ varies. Scaling to $N{=}12$ sharpens the
signature (peak $-0.70$, migrating to $K{=}8$). On mean
time-to-success among successful runs the trade-off produces an
interior optimum at $K{=}8$ ($7.63$ vs.\ $8.04$ independent,
cell-weighted): cell-paired contrasts show $K{=}8$ faster than
$K{=}16$ ($-0.39$, 95\% CI $[-0.80,-0.08]$) and than independent
($-0.40$ $[-0.59,-0.24]$), and statistically tied with
$K{\in}\{2,4\}$. That metric conditions on success. Under a
censoring-aware alternative that charges failed runs the full step
limit, the fast end reverses and independent is marginally best
(paired $\Delta(K{=}8-\text{indep}) = +0.32$ $[+0.05,+0.63]$): fast
sharing's completion collapse costs more successes than its
materialization gain buys back. The interior optimum is therefore a
statement about conditional speed, not unconditional efficiency;
consolidation rate remains a tunable design parameter, and closing
the unconditional gap is what the reservation protocol below
achieves. Oracle interventions confirm the attribution: at fast
cadence, oracle-R/M/C progressively remove the excess time, placing
it at the M$\leftrightarrow$C transition.

\begin{figure}[t]
  \centering
  \includegraphics[width=0.98\columnwidth]{figures/fig_warp_strata.pdf}
  \caption{Provenance stratification of $P(C^\star|M^\star)$ vs.\
  cadence $K$ (peer, scarcity cells, episode-cluster bootstrap CIs).
  The share of locks backed by foreign evidence (warp share, grey)
  grows as sharing accelerates, while that stratum's completion
  (orange) stays below $0.05$ and the own-evidence stratum (blue)
  stays flat: the C-collapse of fast sharing is concentrated in the
  foreign-evidence stratum.}
  \label{fig:strata}
\end{figure}

\textbf{Provenance stratification.} We annotate each $M^\star$ with
the foreign share $\varphi$ of its supporting evidence mass (computed
from the merge's own weights, frozen at lock time) and call locks with
$\varphi{\ge}0.8$ \emph{warps}. The warp share falls from $0.55$ at
$K{=}1$ to $0.11$ at $K{=}8$, while own-evidence locks complete at a
flat $0.26$--$0.34$ and warp locks stay below $0.05$ throughout. The
completion collapse of fast sharing is therefore concentrated in the
foreign-evidence stratum: cadence controls how much of the
collective's commitment relies on peers' evidence, and this stratum
loses the occupancy race. The stratification is not an artefact of
the $0.8$ cutoff: recomputed at thresholds
$\varphi{\in}\{0.5,\dots,0.95\}$, the foreign stratum completes at
${\le}0.03$ while the own stratum stays at $0.26$--$0.50$, at every
threshold and cadence. The two coupling channels also separate
interventionally with arms already in the design: the independent arm
removes peer evidence while leaving target contention in place, and a
decentralized reservation-and-rollback protocol (companion
manuscript) keeps peer evidence while managing contention --- under
it, unconditional success at fast cadence exceeds the $K{=}8$
baseline ($0.614$--$0.620$ vs.\ $0.583$), identifying occupancy
racing rather than evidence quality as the cost channel of fast
sharing. The same stratification curve reproduces on a continuous
VMAS bridge ($0.72\!\to\!0.00$) and on an LLM text substrate
($0.42\!\to\!0.08$).

\textbf{Learned and LLM collectives.} A CommNet-PPO baseline matches
symbolic peer success ($0.667$ vs.\ $0.65$) with structurally
degenerate stages: its real-valued channel is never informatively
empty, so R and M collapse and the stages cannot be read off at
competitive success. A collective of three LLM agents (llama3.1)
sharing symbolic memory reproduces the phenomenology on language.
Paired against a no-communication baseline on identical layouts, raw
sharing lowers success $2/3\!\to\!1/3$ (occupancy racing), and a
decentralized reservation-and-rollback protocol recovers part of it
($0.33\!\to\!0.42$) with rollback latency ${\approx}K$. On a
step-budget sweep with two local models, an opaque success collapse
($1.00\!\to\!0.00$) resolves to pure C-stage failure
($Q^\star{=}R^\star{=}M^\star{=}1$, $C^\star{=}0$) identically for
both models, which identifies the failure as structural rather than
model-specific.

\section{The Contract on Third-Party Agent Stacks}
\label{sec:external}

\begin{figure*}[t]
  \centering
  \includegraphics[width=0.92\textwidth]{figures/fig_external_stages.pdf}
  \caption{Stage profiles of three third-party stacks on four
  engineered variants (llama3.1, $n{=}20$/cell). Q and R saturate
  identically; the structural failures give exact zero columns on
  every stack; the inter-framework spread is confined to M and C.}
  \label{fig:external-stages}
\end{figure*}

The strongest external-validity test of an evaluation protocol is its
application to systems built by others. The contract has two layers:
an \emph{invariant core} --- four event roles (a memory query is
issued; returned evidence satisfies the task; the agent commits to a
target; the task completes) --- and a \emph{system-specific adapter}
that maps logged actions to those roles. This section instruments
three third-party stacks --- an OpenAI-SDK function-calling loop,
LangGraph's ReAct agent, and AutoGen --- with one tool-trace adapter
shared by all three: an episodic household memory behind
\method{recall}/\method{goto}/\method{take} tools; $Q^\star$ on
querying, $R^\star$ on attribution (a returned record names the item
at its true place), $M^\star$ on a commitment (\method{goto}) to the
true place, $C^\star$ on completion within a tool budget. A tool
trace exposes no lock chain, so these are the \emph{marginal event
frequencies} of \S\ref{sec:protocol}: an agent can commit to the true
place without a supporting retrieval, the events are not nested, and
no factorisation is claimed --- the four rates serve as separate
diagnostics. What varies across system classes is the adapter, not
the event roles; two revisions of this adapter are recorded in the
pre-run registration. Four variants engineer one failure each; all
predictions were registered before the runs.

\begin{table}[t]
\caption{Marginal Q/R/M/C event frequencies on third-party stacks (20
episodes/cell; the events are not nested here and no product is
claimed). $M_1$: first commitment targets the true place.}
\label{tab:external}
\centering
\small
\setlength{\tabcolsep}{3.4pt}
\begin{tabular}{llccccc}
\toprule
Stack & Variant & $Q^\star$ & $R^\star$ & $M^\star$ & $C^\star$ & $M_1$ \\
\midrule
OpenAI & control      & 1.00 & 1.00 & .85 & .75 & .25 \\
SDK    & consol.\ gap & 1.00 & \textbf{.00} & .50 & .40 & .25 \\
       & stale ambig. & 1.00 & 1.00 & .90 & .30 & .25 \\
       & tight budget & 1.00 & 1.00 & .25 & \textbf{.00} & .25 \\
\midrule
Lang-  & control      & 1.00 & 1.00 & .70 & .50 & .50 \\
Graph  & consol.\ gap & 1.00 & \textbf{.00} & .20 & .15 & .10 \\
       & stale ambig. & 1.00 & 1.00 & .65 & .60 & .50 \\
       & tight budget & 1.00 & 1.00 & .10 & \textbf{.00} & .10 \\
\midrule
Auto-  & control      & 1.00 & 1.00 & .85 & .55 & .45 \\
Gen    & consol.\ gap & 1.00 & \textbf{.00} & .30 & .20 & .15 \\
       & stale ambig. & 1.00 & 1.00 & .50 & .45 & .35 \\
       & tight budget & 1.00 & 1.00 & .15 & \textbf{.00} & .15 \\
\bottomrule
\end{tabular}
\end{table}

Three findings emerge (Table~\ref{tab:external},
Fig.~\ref{fig:external-stages}). \textbf{(1)~Structural localizations
are exact and framework-independent.} With no consolidated fact,
$R^\star{=}0.00$; with the budget below requirement,
$C^\star{=}0.00$ at $Q^\star{=}R^\star{=}1$. Both hold in $20/20$
episodes on every stack with no per-framework adaptation.
\textbf{(2)~The clean control separates the stacks, and the stages
localize the difference.} Success on identical clean episodes spreads
$0.50$--$0.75$. A leaderboard would report a quality difference and
stop; the stages show a pathology shared by all three stacks
(first-commitment discipline is weak everywhere,
$M_1{=}0.25$--$0.50$: the first \method{goto} frequently ignores the
agent's own just-retrieved evidence) together with different loop
economics: episode-level M recovers to $0.70$--$0.85$ by re-query, and
the stacks differ in how much of that survives to completion ($C/M$
from $0.88$ down to $0.65$). \textbf{(3)~The stale distractor has a
framework-specific cost.} An outdated record captures the first
commitment in $0.20$--$0.40$ of ambiguous episodes, and the completion
cost varies (OpenAI SDK: $0.75\!\to\!0.30$; the others absorb the
detour). Some registered behavioural predictions assumed a clean
control and failed for exactly that reason: the stacks themselves fail
a clean memory task 25--50\% of the time (finding~2). These are
reported as registered failures together with the mechanism.

\textbf{Model robustness (registered before the runs).} Repeating all
240 episodes with a second model (qwen3:4b) confirms the registered
prediction exactly: both structural localizations transfer unchanged
($R^\star{=}0.00$ under the consolidation gap and $C^\star{=}0.00$
under the tight budget, $20/20$ on every stack). The behavioural
profile, by contrast, is a model property that the stages make
explicit. qwen3:4b completes every clean episode ($C^\star{=}1.00$ on
all three stacks) with first-commitment discipline
$M_1{=}0.90$--$0.95$ where llama3.1 managed $0.25$--$0.50$, and is
barely distracted by the stale record ($0.05$--$0.30$). The exception
is AutoGen, which drops qwen's $M_1$ to $0.40$; the decomposition thus
separates the model's contribution from the effect of the framework's
prompting loop. The full second-model table is in the supplement.

\section{Related Work}

Cognitive-map and semantic-map lines motivate explicit place
abstraction \citep{tolman1948cognitive,gupta2017cognitive,
savinov2018semi}; goal-conditioned RL conditions on goals or
embeddings \citep{andrychowicz2017hindsight}; multi-agent RL studies
learned channels \citep{sukhbaatar2016commnet,foerster2018ctde} and
gossip protocols \citep{boyd2006gossip}; LLM-agent systems add
explicit memory layers evaluated end-to-end
\citep{wang2023voyager,shinn2023reflexion,packer2023memgpt,
park2023generative,wu2023autogen}. The protocol is orthogonal to
these lines: it proposes no architecture, makes any of them
stage-diagnosable from logs, and yields a slope-level cadence
prediction testable on those logs.

\textbf{Stage-wise evaluation.} The closest line decomposes
retrieval-augmented generation: RAGAS scores retrieval and generation
quality separately \citep{es2024ragas}, RAGChecker refines this to
claim-level diagnostics meta-evaluated against human judgments
\citep{ru2024ragchecker}, and AgentBoard replaces final success with
analytic progress rates for multi-turn LLM agents
\citep{ma2024agentboard}. Q/R/M/C differs on three axes: its events
are logged actions of a task-executing agent, not LLM-judged text
quality; it models commitment (M) and the multi-agent coupling
channels (shared memory, occupancy) that RAG metrics do not
represent; and its diagnoses are validated interventionally (oracle
stages; engineered structural faults) rather than by correlation with
human ratings. Process mining and conformance checking compare event
logs against stage models
\citep{vanderaalst2016process,carmona2018conformance}; the protocol
shares that log-level discipline and adds agent-specific event
semantics with intervention-based validation. Process reward models
score intermediate reasoning steps to supervise training
\citep{lightman2024verify}; Q/R/M/C is an evaluation-time instrument
for memory pipelines.

\section{Limitations}

Primary substrates are grid worlds; VMAS is a directional continuous
check; the LLM studies use small local models at modest $n$; place
identity across agents assumes a shared coordinate frame (removed
constructively in a companion manuscript); Empirical Claim~1 remains
slope-level; a derivation under a generative model would make it a
theorem. The four stages are one resolution level: sub-stage causes
(accumulation within R; ranking, grounding, and staleness handling
within M) require within-link forensics, which the protocol scopes
but does not automate. Full tables, proofs of the estimability arguments,
assumption stress tests, effect-size methodology, scale-up, and all
registered-prediction files are in the supplement.

\section{Conclusion}

Q/R/M/C replaces the scalar verdict ``the agent's memory failed''
with a diagnosis that names a stage. In our studies this located a
consolidation limit in a single agent, a sharing-rate trade-off priced
by occupancy in a collective, and a commitment-discipline gap in
off-the-shelf agent stacks; in several cases the diagnosis was
specific enough to act on. The protocol requires a logging contract
rather than architectural modification, and the resulting rates are
directly comparable across systems that expose a query/lock
interface; for opaque learned policies the protocol measures stage
observability itself.

\bibliography{aaai27refs}

% Reproducibility checklist (pages beyond the 7-page technical limit
% are permitted for references and the checklist).
\def\isChecklistMainFile{}
\input{ReproducibilityChecklist_qrmc_aaai27.tex}

\end{document}
